% fig_multilevel_metagraph.tex — Вертикальные колонки без пересечений
\documentclass[tikz, border=3mm]{standalone}
\usepackage{fontspec}
\usepackage[russian]{babel}
\setsansfont{PT Sans}
\usetikzlibrary{positioning, arrows.meta, backgrounds, fit}
\begin{document}
\begin{tikzpicture}[
    levelbox/.style={rectangle, draw=black, thick, rounded corners, fill=#1, inner sep=6pt, font=\small\sffamily, minimum height=2.0cm},
    node/.style={rectangle, draw=black, thick, fill=white, font=\scriptsize\sffamily, inner sep=3pt, minimum width=2.2cm, minimum height=0.2cm,align=center},
    link/.style={-{Stealth[length=4pt,width=6pt]}, thick, black!70},
    label/.style={font=\scriptsize\sffamily}
]

% === Основные параметры ===
\def\colsep{5.5cm}   % расстояние между колонками
\def\rowsep{1.5cm}   % расстояние между узлами в колонке

% =============== КОЛОНКА 1: Целостность ===============
% Стратегический уровень
\node[node, fill=blue!20] (prop1) at (0, 0) {Нарушение\\целостности};
\node[node, fill=blue!30, below=\rowsep of prop1] (goal1) {Изменение\\финансовых\\транзакций};
\node[node, fill=orange!20, below=\rowsep of goal1] (t1) {T1102\\Web\\Service};
\node[node, fill=green!20, below=\rowsep of t1] (m1a) {powershell.exe};
\node[node, fill=green!20, below=\rowsep of m1a] (m1b) {wscript.exe};
\node[node, fill=green!20, below=\rowsep of m1b] (m1c) {SQL\_inject-PoC.py};

% Связи
\draw[link] (prop1) -- (goal1);
\draw[link] (goal1) -- node[label, right] {реализует} (t1);
\draw[link] (t1) -- (m1a);
\draw[link] (m1a) -- (m1b);
\draw[link] (m1b) -- (m1c);

% =============== КОЛОНКА 2: Конфиденциальность ===============
\node[node, fill=blue!20] (prop2) at (\colsep, 0) {Нарушение\\конфиденциальности};
\node[node, fill=blue!30, below=\rowsep of prop2] (goal2) {Хищение\\персональных\\данных};
\node[node, fill=orange!20, below=\rowsep of goal2] (t2) {T1566\\Phishing};
\node[node, fill=green!20, below=\rowsep of t2] (m2a) {Outlook.exe};
\node[node, fill=green!20, below=\rowsep of m2a] (m2b) {acrobat.exe};
\node[node, fill=green!20, below=\rowsep of m2b] (m2c) {malware.exe};

\draw[link] (prop2) -- (goal2);
\draw[link] (goal2) -- (t2);
\draw[link] (t2) -- (m2a);
\draw[link] (m2a) -- (m2b);
\draw[link] (m2b) -- (m2c);

% =============== КОЛОНКА 3: Доступность ===============
\node[node, fill=blue!20] (prop3) at (2*\colsep, 0) {Нарушение\\доступности};
\node[node, fill=blue!30, below=\rowsep of prop3] (goal3) {DDoS-атака\\на сервис};
\node[node, fill=orange!20, below=\rowsep of goal3] (t3) {T1499\\Endpoint\\DoS};
\node[node, fill=green!20, below=\rowsep of t3] (m3a) {LOIC};
\node[node, fill=green!20, below=\rowsep of m3a] (m3b) {hping3};
\node[node, fill=green!20, below=\rowsep of m3b] (m3c) {slowloris.py};

\draw[link] (prop3) -- (goal3);
\draw[link] (goal3) -- (t3);
\draw[link] (t3) -- (m3a);
\draw[link] (m3a) -- (m3b);
\draw[link] (m3b) -- (m3c);

% === Фоны уровней с подписями внутри ===
\begin{scope}[on background layer]
    % Стратегический
    \node[levelbox=blue!10, fit=(prop1) (prop3), inner ysep=8pt] (bg1) {};
    \node[font=\bfseries\sffamily, above=2pt of bg1.north, anchor=south] {Стратегический уровень};
    
    % Тактический
    \node[levelbox=orange!10, fit=(t1) (t3), inner ysep=8pt] (bg2) {};
    \node[font=\bfseries\sffamily, above=2pt of bg2.north, anchor=south] {Тактический уровень (MITRE ATT\&CK)};
    
    % Методический
    \node[levelbox=green!10, fit=(m1c) (m3c), inner ysep=8pt] (bg3) {};
    \node[font=\bfseries\sffamily, above=2pt of bg3.north, anchor=south] {Методический уровень};
\end{scope}

\end{tikzpicture}
\end{document}